\documentclass[11pt]{amsart}

\usepackage[T1]{fontenc}
\usepackage{lmodern}
\usepackage{microtype}
\usepackage{amsmath,amssymb,amsthm}
\usepackage[hidelinks]{hyperref}

\hypersetup{
  pdftitle={Counterexamples to Gui's Sorting Conjecture for Reduced Kronecker Coefficients}
}

\newtheorem{theorem}{Theorem}
\newtheorem{lemma}[theorem]{Lemma}
\newtheorem{corollary}[theorem]{Corollary}

\newcommand{\sort}{\operatorname{sort}}

\title[Gui's sorting conjecture]
{Counterexamples to Gui's Sorting Conjecture for Reduced Kronecker Coefficients}

\date{}

\subjclass[2020]{Primary 20C30; Secondary 05E10}
\keywords{reduced Kronecker coefficient, stable Kronecker coefficient,
sorting, counterexample}

\begin{document}

\begin{abstract}
Gui conjectured that splitting the decreasing rearrangement of the combined
parts of two partitions into its odd- and even-indexed subsequences produces a
coefficientwise larger stable tensor product.  We disprove this conjecture with
the infinite family
\[
  \overline g_{(k,1),(k-1)}^{(1)}=0<1=
  \overline g_{(k),(k-1,1)}^{(1)}
  \qquad(k\ge2).
\]
The case $k=2$ is, up to interchanging the factors, the unique counterexample
among pairs of nonempty partitions with at most four boxes in total.  The same
family also disproves the multipartition extension stated as Conjecture~5.7 in
the published version of Gui's paper.
\end{abstract}

\maketitle

Let $g_{\lambda,\mu}^{\nu}$ denote the ordinary Kronecker coefficient.  For a
partition $\alpha$ and all sufficiently large $N$, set
\[
  \alpha[N]=(N-|\alpha|,\alpha).
\]
Murnaghan's stability theorem defines the reduced Kronecker coefficient by
\[
  \overline g_{\alpha,\beta}^{\gamma}
  =\lim_{N\to\infty}g_{\alpha[N],\beta[N]}^{\gamma[N]};
\]
see \cite{Murnaghan}.

For partitions $\lambda$ and $\mu$, let
$(\delta_1,\delta_2,\ldots)$ be the decreasing rearrangement of all their
parts, and define
\[
  \sort_1(\lambda,\mu)=(\delta_1,\delta_3,\ldots),
  \qquad
  \sort_2(\lambda,\mu)=(\delta_2,\delta_4,\ldots).
\]
Gui's Conjecture~5.5 \cite{Gui} asserts that
\begin{equation}\label{eq:gui}
  \overline g_{\sort_1(\lambda,\mu),\sort_2(\lambda,\mu)}^{\nu}
  \ge \overline g_{\lambda,\mu}^{\nu}
  \qquad\text{for every partition }\nu.
\end{equation}

For a partition $\alpha$, let $R(\alpha)$ be the set of partitions obtained by
deleting one removable corner.  For a statement $P$, write $\mathbf 1[P]=1$
if $P$ holds and $\mathbf 1[P]=0$ otherwise.

\begin{lemma}\label{lem:one}
For all partitions $\alpha$ and $\beta$,
\[
  \overline g_{\alpha,\beta}^{(1)}
  =|R(\alpha)\cap R(\beta)|
   +\mathbf 1[\alpha\in R(\beta)]
   +\mathbf 1[\beta\in R(\alpha)].
\]
\end{lemma}

\begin{proof}
For a partition $\sigma\vdash N$, write $[\sigma]$ for the corresponding
irreducible $S_N$-module.  The permutation module on $N$ letters is
$[N]\oplus[N-1,1]$, and the tensor identity gives
\[
  [\sigma]\otimes\mathbb C^N
  \cong
  \operatorname{Ind}_{S_{N-1}}^{S_N}
  \operatorname{Res}_{S_{N-1}}^{S_N}[\sigma].
\]
The branching rule therefore yields, for $N\ge2$ and
$\sigma,\rho\vdash N$,
\begin{equation}\label{eq:branching}
  g_{\sigma,(N-1,1)}^{\rho}
  =|R(\sigma)\cap R(\rho)|-\delta_{\sigma,\rho}.
\end{equation}
For all sufficiently large $N$,
\[
  R(\alpha[N])
  =\{\alpha[N-1]\}\cup
    \{\tau[N-1]:\tau\in R(\alpha)\},
\]
and likewise for $\beta$.  By symmetry of ordinary Kronecker coefficients,
\eqref{eq:branching} with $\sigma=\alpha[N]$ and $\rho=\beta[N]$ computes
$g_{\alpha[N],\beta[N]}^{(N-1,1)}$.  The two first-row removals coincide
exactly when $\alpha=\beta$, and this contribution is cancelled by the
Kronecker delta.  The lower-row removals contribute
$|R(\alpha)\cap R(\beta)|$, while the two cross terms occur exactly when
$\alpha\in R(\beta)$ or $\beta\in R(\alpha)$.  Taking the stable value proves
the formula.
\end{proof}

\begin{theorem}\label{thm:main}
For every integer $k\ge2$,
\[
  \overline g_{(k,1),(k-1)}^{(1)}=0,
  \qquad
  \overline g_{(k),(k-1,1)}^{(1)}=1.
\]
Consequently, Gui's Conjecture~5.5 is false.
\end{theorem}

\begin{proof}
The partitions $(k)$ and $(k-1,1)$ have the common one-box predecessor
$(k-1)$ and no other common predecessor.  Hence Lemma~\ref{lem:one} gives
\[
  \overline g_{(k),(k-1,1)}^{(1)}=1.
\]
The sizes of $(k,1)$ and $(k-1)$ differ by two, so every term on the
right-hand side of Lemma~\ref{lem:one} vanishes.  Thus
\[
  \overline g_{(k,1),(k-1)}^{(1)}=0.
\]
Finally, the decreasing rearrangement of the parts of $(k)$ and $(k-1,1)$ is
$(k,k-1,1)$, and therefore
\[
  \sort_1\bigl((k),(k-1,1)\bigr)=(k,1),
  \qquad
  \sort_2\bigl((k),(k-1,1)\bigr)=(k-1).
\]
At $\nu=(1)$, inequality \eqref{eq:gui} would therefore require $0\ge1$.
\end{proof}

\begin{corollary}\label{cor:minimal}
Up to interchanging the factors, $((2),(1,1))$ is the unique pair of nonempty
partitions $(\lambda,\mu)$ with $|\lambda|+|\mu|\le4$ for which
\eqref{eq:gui} fails for some partition $\nu$.
\end{corollary}

\begin{proof}
For total size at most three, sorting fixes the pair up to interchange, so
\eqref{eq:gui} holds with equality.  At total size four, the only unordered
pairs not fixed by sorting are
\[
  ((1),(1,1,1))
  \qquad\text{and}\qquad
  ((2),(1,1)).
\]
Gui's Theorem~5.6 \cite{Gui} proves \eqref{eq:gui} for the first pair, since
both factors are columns.  The second pair fails by Theorem~\ref{thm:main}
with $k=2$.
\end{proof}

\begin{corollary}
The multipartition statement published as Gui's Conjecture~5.7 is false.
\end{corollary}

\begin{proof}
For two input partitions, the two partitions denoted
$(\lambda\cup\mu)^{[1,2]}$ and $(\lambda\cup\mu)^{[2,2]}$ in
\cite[Conjecture~5.7]{Gui} are precisely
$\sort_1(\lambda,\mu)$ and $\sort_2(\lambda,\mu)$.  Thus the two-factor case
of that statement is Conjecture~5.5.
\end{proof}

\begin{thebibliography}{9}

\bibitem{Gui}
T.~Gui,
\emph{Conjectures on the reduced Kronecker coefficients},
S\'em. Lothar. Combin. \textbf{93B} (2025), Art.~1, 12~pp.,
\href{https://www.mat.univie.ac.at/~slc/wpapers/FPSAC2025/1.pdf}
{published version}.

\bibitem{Murnaghan}
F.~D.~Murnaghan,
\emph{The analysis of the Kronecker product of irreducible representations
of the symmetric group},
Amer. J. Math. \textbf{60} (1938), no.~3, 761--784,
\href{https://doi.org/10.2307/2371610}{doi:10.2307/2371610}.

\end{thebibliography}

\end{document}